Classical nova V339 Del (Nova Delphini 2013) was discovered by Koichi Itagaki
at 6.8 magnitude on 14 August 2013 at 14:01 UT ($\mathrm{MJD} = 56518.584$).
Both professional and amateur astronomers analysed the photometry and
spectroscopy of the nova. Less than 10 hours after the alert, when the night
falls at the Piszk\'estet\H{o} Mountain Station of the Konkoly Observatory of
the Hungarian Academy of Sciences, Hungarian astronomers took the first
spectrum data of the nova using the eShel echelle spectrograph in the Gothard
Astrophysical Observatory of Lor\'and E\"otv\"os University attached to the
1 meter telescope of the Konkoly Observatory.

Refer to Fig.~1.1 and Fig.~1.2 to complete the questions. The larger versions
of Fig.~1.1 and Fig.~1.2 are found on separate A3 papers.

Fig.~1.1 shows the visual light curve of the nova based on the data
downloaded from the website of AAVSO (American Association of Variable Star
Observers). On the horizontal and vertical axes, the Modified Julian Date
($\mathrm{MJD} = \mathrm{JD} - 2\,400\,000.5$) of the observations and the
Johnson $V$ magnitudes are plotted, respectively. The grey circles (about
38000 data points) represent the measured values, while the continuous black
line is the result of smoothing the data with a Gaussian filter (Full Width
at Half Maximum $= 0.5$\,day) to define an ``average'' light curve from the
data points.

The rate of decline can be characterized by the values $t_2$ and $t_3$, which
show the time interval in days in which a nova fades from its maximum
brightness by 2 and 3 magnitudes.

A few empirical formulae between the peak of the absolute magnitude in the
$V$ band ($M_0$) and $t_2$, $t_3$ values can be found in the following
literature:
\begin{description}
\item[(a)] $M_0 = -7.92 - 0.81\arctan\dfrac{1.32 - \log t_2}{0.23}$ \\
  (Della Valle, M. \& Livio, M.: 1995, ApJ 452, 704)
\item[(b)] $M_0 = -11.32 + 2.55\log t_2$ \\
  (Downes, R.A. \& Durbeck, H.W.: 2000, AJ 120, 2007)
\item[(c)] $M_0 = -11.99 + 2.54\log t_3$ \\
  (Downes, R.A. \& Durbeck, H.W.: 2000, AJ 120, 2007)
\end{description}

The $E(B-V)$ color excess of Nova Del 2013 (Chochol, D. et al.: 2014,
Contrib. Astron. Obs. Skalnat\'e Pleso 43, 330) is:
\[ E(B-V) = 0.184 \pm 0.035 \]

Fig.~1.2 shows the nova spectra taken in the wavelength region around the
H$\alpha$ line on six consecutive nights before and after the time of the
maximum brightness ($t_0$). The individual spectra have been shifted
vertically for clarity. The Modified Julian Dates (MJD) of the observations
are listed on the right hand side of each spectrum slice.

The H$\alpha$ line shows the so-called P Cygni profile with very broad wings,
which is typical not of novae only, but is present in almost all spectral
types and is a reliable sign of a massive radial motion of matter ejected
from the star. The P Cygni profile is composed of a strong, broad emission
peak which is considered to be centered at the rest wavelength in air
$\lambda_0$ of the line -- for H$\alpha$ $\lambda_0 = 6562.82$\,\AA{} -- and
a usually weaker, blueshifted absorption component. The expansion (radial)
velocity of the shell can be approximated from the measured wavelength
$\lambda$ of the absorption peak using the well known Doppler formula
connecting the displacement $\Delta\lambda = \lambda - \lambda_0$, the
radial velocity $v_{\mathrm{r}}$, and $c$, the speed of light.

Assume that the H$\alpha$ line showing P Cygni profile is excited in the
outermost part of the spherically expanding shell, and its extent at the
moment of taking the first spectrum was still negligible.

\begin{description}
\item[(a)] From Fig.~1.1, estimate the Modified Julian Date of the peak
  magnitude ($\mathrm{MJD}_0$) and the value of the peak magnitude itself.
  Consider the error of this brightness value to be $0.05^{\mathrm{m}}$.
  \hfill[3\,p]
\item[(b)] Estimate the Modified Julian Dates based on the time interval
  (days) in which the nova has faded by 2 and 3 magnitudes, then calculate
  $t_2$ and $t_3$ values. \hfill[6\,p]
\item[(c)] With reference to $t_2$ and $t_3$ from (b), determine the peak
  absolute magnitude of the nova using all three empirical formulae listed
  earlier, calculate their mean ($M_0$) and their standard deviation, and
  consider this latter as the uncertainty of $M_0$. \hfill[5\,p] \\
  The formula for standard deviation:
  \[ \sigma = \sqrt{\frac{\sum_{i=1}^{n}(x_i - \bar{x})^2}{n-1}} \]
\item[(d)] Using the value of the color excess $E(B-V)$, determine the
  interstellar extinction $A_V$ and its uncertainty in the direction of the
  nova. Use $R_V = 3.1$, without error. \hfill[4\,p]
\item[(e)] Estimate the distance to the nova and its uncertainty. Give the
  result in kpc. \hfill[11\,p]
\item[(f)] Measure the central wavelengths of the P Cygni absorption
  features plotted in Fig.~1.2 (refer to magnified version), and calculate
  the corresponding radial velocities. No error estimation is needed.
  \hfill[14\,p]
\item[(g)] Plot these radial velocities against the Modified Julian Dates of
  the observations. \hfill[6\,p]
\item[(h)] From the graph in (g), estimate the physical radius of the
  envelope at the end of the time interval. Give the answer in astronomical
  units (au). \hfill[7\,p]
\item[(i)] Knowing the distance to the nova and the physical radius of the
  spherical envelope 5 days after the discovery, estimate the apparent
  angular diameter of the envelope then. \hfill[4\,p]
\end{description}
